% ============================================================
% VII. DISCUSSION
% ============================================================
\section{Discussion}
\label{sec:discussion}

% ---- A. Practical Trade-offs ----
\subsection{When to Use Each Mode}

The three approximation modes serve different use cases in the robotics pipeline.
\textbf{Convex mode} produces the geometrically tightest approximation and is the best choice when the target simulator or collision checker handles convex meshes natively.
It requires no parametric decomposition and preserves the mesh's convex envelope faithfully.
However, it remains a mesh representation and does not admit the constant-time distance evaluation of analytic primitives.

\textbf{Sphere mode} is well-suited to pipelines that exploit sphere algebra---GPU-based collision pipelines, signed-distance fields via sphere-sums, or planners that can use spherical bounding volume hierarchies.
The medial tree produces multiple small spheres that can nestle into concave regions, often giving tighter overall coverage than capsules on shapes with deep concavities.
The trade-off is higher primitive count: a typical sphere-tree approximation requires many more primitives than the equivalent capsule fit.

\textbf{Capsule mode} is the best choice when smooth, low-count analytic primitives are desired.
Capsules provide constant-time point-to-capsule distance, making them attractive for collision-atlas precomputation, swept-volume computation along trajectories, and analytical signed-distance fields for motion planning.
These analytic primitives support constant-time collision distance evaluation, making them suitable for collision-checking in motion planning (e.g., via FCL~\cite{pan2012fcl} or OMPL), real-time simulation, and differentiable collision pipelines.
The principal-axis sectioning assumption means that capsules are most effective on elongated, roughly tubular link geometries---precisely the shape of most robot arm links.

% ---- B. URDF Compatibility ----
\subsection{URDF Compatibility Constraint}

A significant practical constraint is the absence of a native capsule element in the URDF specification.
Our toolchain must decompose each capsule into a cylinder plus two spheres for the collision URDF, which introduces a subtle issue: a simulator loading the URDF reconstructs the collision shape as three separate collision primitives rather than as a single capsule.
This means the simulator does not benefit from the capsule's analytic distance function.
The JSON sidecar addresses this limitation by storing the canonical capsule parameters separately, but bridging the gap between URDF-compatible emission and analytic usability requires downstream support for reading the sidecar.

The broader implication is that the robotics community would benefit from extending URDF (or its successor formats) to include a standard capsule primitive.
Until such an extension is adopted, toolchains like ours must maintain dual outputs.

% ---- C. Coverage vs. Tightness ----
\subsection{Coverage--Tightness Trade-off and Limitations}

Conservative coverage---requiring every mesh vertex to lie inside at least one capsule---creates inherent tension with tightness. Our splitting strategy, guided by $r/\mathrm{binMed}$ and $\mathrm{capV/aabb}$, selectively splits capsules where radius inflation indicates geometric inefficiency, avoiding over-splitting in tight regions.

The principal-axis sectioning approach assumes each link has a clear longest axis for capsule orientation. While this holds for standard manipulator links, it may perform poorly on branched, flat, or spherical geometry. Additionally, PCA captures only global elongation; locally curved links may benefit from curved capsule chains or articulated capsule hierarchies.

% ---- E. Reproducibility ----
\subsection{Reproducibility}

The toolchain ensures reproducible execution through deterministic random seeds, fixed presets, and a containerized pipeline. Validation gates accept user-defined quality criteria (e.g., $\mathrm{capV/aabb} < 2.5$) for automated pass/fail checking.
